%============================================================
% CH2 · P4 · PROBLEM 1 — ALL PATHS FROM SOURCE TO TARGET (LeetCode 797)
%============================================================
\subsection{Problem 2.4.1: All Paths From Source to Target}
\problemlabel{2.4.1}{ssec:ch2-all-paths}

\problemcard{2.4.1}{All Paths From Source to Target}{LeetCode 797}{https://leetcode.com/problems/all-paths-from-source-to-target/}{\exeasy\ \lc{Medium}\ (the first backtracking program; the output itself is exponential)}{Backtracking with a path stack; in a DAG no \texttt{visited} array is needed at all; alternative: memoised path lists per vertex}

\subsubsection*{1. Problem statement}

A directed acyclic graph (DAG) has $n$ vertices $0..n-1$; \texttt{graph[i]} lists the vertices $j$ with an edge $i \to j$. Return \textbf{all} paths from vertex $0$ to vertex $n-1$, in any order.

\textbf{Constraints:} $2 \le n \le 15$; $0 \le \texttt{graph[i][j]} < n$, $\texttt{graph[i][j]} \ne i$; the entries of one list are distinct; the graph is guaranteed to be a DAG.

\begin{table}[H]
\centering\small
\begin{tabular}{@{}lll@{}}
\toprule
\textbf{graph} & \textbf{Official output} & \textbf{Our output (backtracking)} \\
\midrule
\texttt{[[1,2],[3],[3],[]]} & \texttt{[[0,1,3],[0,2,3]]} & \texttt{[[0,1,3],[0,2,3]]} \\
\texttt{[[4,3,1],[3,2,4],[3],[4],[]]} & \texttt{[[0,4],[0,3,4],[0,1,3,4],} & identical, same order \\
 & \texttt{\ [0,1,2,3,4],[0,1,4]]} & \\
\bottomrule
\end{tabular}
\end{table}

\begin{diagrambox}
\begin{tikzpicture}[x=1.15cm, y=1cm]
    \begin{scope}
        \node[gvsrc] (a0) at (0,0) {0}; \node[gv] (a1) at (1.5,0.8) {1};
        \node[gv] (a2) at (1.5,-0.8) {2}; \node[gvdst] (a3) at (3,0) {3};
        \foreach \a/\b in {0/1, 0/2, 1/3, 2/3} \draw[gde] (a\a) -- (a\b);
        \node[font=\scriptsize\bfseries] at (1.5,-1.75) {Example 1: 2 paths};
    \end{scope}
    \begin{scope}[xshift=5.6cm]
        \node[gvsrc] (b0) at (0,0) {0}; \node[gv] (b1) at (1.4,1.2) {1}; \node[gv] (b2) at (2.9,1.2) {2};
        \node[gv] (b3) at (2.9,-0.6) {3}; \node[gvdst] (b4) at (4.3,0) {4};
        \draw[gde] (b0) -- (b1); \draw[gde] (b0) -- (b3); \draw[gde] (b0) to[bend right=50] (b4);
        \draw[gde] (b1) -- (b2); \draw[gde] (b1) -- (b3); \draw[gde] (b1) -- (b4);
        \draw[gde] (b2) -- (b3); \draw[gde] (b3) -- (b4);
        \node[font=\scriptsize\bfseries] at (2.2,-1.75) {Example 2: 5 paths};
    \end{scope}
\end{tikzpicture}
\figcap{The two examples. In Example 2 vertex 3 is reached along three different paths ($0 3$, $0 1 3$, $0 1 2 3$), vertex 4 along five.}
\end{diagrambox}

\subsubsection*{2. The graph model}

\begin{modelbox}[Graph Model --- All Paths From Source to Target]
\renewcommand{\arraystretch}{1.35}
\begin{tabularx}{\linewidth}{@{}>{\bfseries}l X@{}}
Vertices & the $n \le 15$ given vertices \\
Edges & directed, given as adjacency lists; acyclic; no self-loops, no parallel edges \\
Weights & none \\
Question & list every path $0 \leadsto n-1$; in a DAG every path is automatically simple \\
Size & at most $\binom{15}{2} = 105$ edges --- the graph is tiny; the \textbf{output} is the problem: up to $2^{n-2} = 8192$ paths with $69\,632$ numbers in total \\
\end{tabularx}
\end{modelbox}

\subsubsection*{3. Key observations}

\begin{itemize}[leftmargin=*]
    \item \textbf{No marks needed.} A path in a DAG can never return to a vertex it has used: that would close a cycle. So the \texttt{onPath} test of template 1 is always false and can be dropped; only the path stack is needed.
    \item \textbf{A path is its set of middle vertices.} Fix a topological order. The vertices of any path appear in that order, so a path is determined by \emph{which} of the $n-2$ middle vertices it uses. Hence there are at most $2^{n-2}$ paths --- and the \textbf{complete DAG} (an edge $i \to j$ for all $i < j$) has exactly that many, because every subset in increasing order is a path.
    \item \textbf{Output-sensitive.} The total output size is at most $2^{13} \cdot 15$ numbers; any correct algorithm must spend at least that. The goal is to spend not much more.
    \item \textbf{Dead ends.} A branch into vertices that cannot reach $n-1$ produces no output but costs time; a reverse search from $n-1$ (Pattern~5) finds them in advance.
\end{itemize}

\begin{diagrambox}
\begin{tikzpicture}[x=1.25cm, y=1cm]
    \foreach \i in {0,...,4} \node[gv] (c\i) at (\i,0) {\i};
    \node[gvsrc] at (c0) {0}; \node[gvdst] at (c4) {4};
    \foreach \i/\j in {0/1, 1/2, 2/3, 3/4} \draw[gde] (c\i) -- (c\j);
    \foreach \i/\j/\b in {0/2/30, 1/3/30, 2/4/30} \draw[gde] (c\i) to[bend left=\b] (c\j);
    \foreach \i/\j/\b in {0/3/38, 1/4/38} \draw[gde] (c\i) to[bend right=\b] (c\j);
    \draw[gde] (c0) to[bend right=55] (c4);
    \node[font=\scriptsize\bfseries] at (2,-2.3) {complete DAG, $n = 5$};
    \begin{scope}[xshift=7.2cm]
        \node[font=\scriptsize\bfseries] at (1.5,1.05) {middle set $\to$ path};
        \foreach \s/\p/\y in {{\{\}}/0 4/0.4, {\{1\}}/0 1 4/-0.05, {\{2\}}/0 2 4/-0.5, {\{3\}}/0 3 4/-0.95,
                              {\{1,2\}}/0 1 2 4/-1.4, {\{1,3\}}/0 1 3 4/-1.85, {\{2,3\}}/0 2 3 4/-2.3, {\{1,2,3\}}/0 1 2 3 4/-2.75}{
            \node[font=\scriptsize, anchor=east] at (0.9,\y) {$\s$};
            \node[font=\scriptsize\ttfamily, anchor=west] at (1.2,\y) {\p};
        }
        \draw[->, darkgray] (0.95,-1.2) -- (1.15,-1.2);
    \end{scope}
\end{tikzpicture}
\figcap{Why $2^{n-2}$ is the maximum and is attained: in the complete DAG each of the $2^3 = 8$ subsets of $\{1,2,3\}$, read in increasing order, is a path from 0 to 4.}
\end{diagrambox}

\subsubsection*{4. Approaches}

\begin{bruteforcebox}
\textbf{Idea: try every candidate set.} By the second observation it is enough to try every subset of the middle vertices: sort it in a topological order (computed by Kahn's algorithm, Chapter~5), put 0 in front and $n-1$ at the end, and keep it if every consecutive pair is an edge.

\begin{lstlisting}[style=cpp]
#include <bits/stdc++.h>
using namespace std;

// Brute force: every subset of the middle vertices, in topological order, checked edge by edge.
class Solution {
public:
    vector<vector<int>> allPathsSourceTarget(vector<vector<int>>& graph) {
        int n = graph.size();
        vector<int> indeg(n, 0), order;                // Kahn's algorithm for a topological order
        for (auto& l : graph) for (int v : l) indeg[v]++;
        queue<int> q;
        for (int v = 0; v < n; v++) if (indeg[v] == 0) q.push(v);
        while (!q.empty()) {
            int u = q.front(); q.pop();
            order.push_back(u);
            for (int v : graph[u]) if (--indeg[v] == 0) q.push(v);
        }
        vector<vector<char>> edge(n, vector<char>(n, 0));
        for (int u = 0; u < n; u++) for (int v : graph[u]) edge[u][v] = 1;
        vector<int> mid;
        for (int v : order) if (v != 0 && v != n - 1) mid.push_back(v);
        int k = mid.size();
        vector<vector<int>> res;
        for (int mask = 0; mask < (1 << k); mask++) {
            vector<int> p(1, 0);
            for (int i = 0; i < k; i++) if (mask >> i & 1) p.push_back(mid[i]);
            p.push_back(n - 1);
            bool ok = true;
            for (size_t i = 0; i + 1 < p.size() && ok; i++) ok = edge[p[i]][p[i + 1]];
            if (ok) res.push_back(p);
        }
        return res;
    }
};

int main() {
    vector<vector<int>> g = {{4, 3, 1}, {3, 2, 4}, {3}, {4}, {}};
    for (auto& p : Solution().allPathsSourceTarget(g)) {
        for (int x : p) cout << x << " ";
        cout << "/ ";
    }
    cout << "\n";                                      // 0 4 / 0 1 4 / 0 3 4 / 0 1 3 4 / 0 1 2 3 4 /
}
\end{lstlisting}

\begin{complexitybox}
\textbf{Time:} $2^{n-2}$ subsets, \bigO{n} each: $8192 \cdot 15 \approx 1.2 \cdot 10^5$ steps for $n = 15$ --- fast, because $n$ is so small. But it pays $2^{n-2}$ even when the graph has a single path. \textbf{Space:} \bigO{n^2} plus the output. \textbf{Verdict:} \ac{}-fast for $n \le 15$ and an ideal reference for testing; useless as soon as $n$ grows.
\end{complexitybox}
\end{bruteforcebox}

\begin{algorithmbox}[Approach 2: backtracking with a path stack]
\textbf{Idea.} Template 1 of the pattern without the marks: push $u$, and if $u = n-1$ record a copy of the path, otherwise recurse into every successor; pop $u$ on exit.

\begin{lstlisting}[style=cpp]
#include <bits/stdc++.h>
using namespace std;

// All Paths From Source to Target (LeetCode 797): backtracking. In a DAG a path can never
// revisit a vertex, so no visited/onPath array is needed -- only the path stack.
class Solution {
    vector<vector<int>> res;
    vector<int> path;
    void dfs(const vector<vector<int>>& g, int u) {
        path.push_back(u);                             // enter u
        if (u == (int)g.size() - 1) res.push_back(path);
        else
            for (int v : g[u]) dfs(g, v);
        path.pop_back();                               // exit u
    }
public:
    vector<vector<int>> allPathsSourceTarget(vector<vector<int>>& graph) {
        res.clear();
        dfs(graph, 0);
        return res;
    }
};

int main() {
    vector<vector<int>> g1 = {{1, 2}, {3}, {3}, {}};
    vector<vector<int>> g2 = {{4, 3, 1}, {3, 2, 4}, {3}, {4}, {}};
    for (auto* g : {&g1, &g2}) {
        for (auto& p : Solution().allPathsSourceTarget(*g)) {
            for (int x : p) cout << x << " ";
            cout << "/ ";
        }
        cout << "\n";
    }
    // 0 1 3 / 0 2 3 /
    // 0 4 / 0 3 4 / 0 1 3 4 / 0 1 2 3 4 / 0 1 4 /
}
\end{lstlisting}

\begin{complexitybox}
\textbf{Time:} one call per path from 0 to \emph{any} vertex, plus a copy of every finished path. On the complete DAG with $n = 15$: $16\,384$ calls ($= 2^{14}$: there are $2^{v-1}$ paths from 0 to each $v \ge 1$, plus the root) and $8192$ copies of total length $69\,632$, in 3\,ms. \textbf{Space:} \bigO{n} for the stack and the path, plus the output. \textbf{Verdict:} \ac{}; within a factor 2 of the output size on the worst case.
\end{complexitybox}
\end{algorithmbox}

\begin{invariantbox}[Why it is correct]
\textbf{Invariant.} When \texttt{dfs(u)} starts, \texttt{path} holds a path $0 = x_0 \to x_1 \to \dots \to x_k = u$ of the graph; when it returns, \texttt{path} is $x_0, \dots, x_{k-1}$ again. (Push and pop are paired; induction over the calls, exactly as in the pattern's invariant.)

\emph{Soundness.} Every recorded \texttt{path} ends in $n-1$ and consists of edges, so it is a path $0 \leadsto n-1$; in a DAG it is automatically simple.

\emph{Completeness and no duplicates.} A path $0 = y_0 \to \dots \to y_m = n-1$ is produced by the chain of calls $\texttt{dfs}(y_0), \dots, \texttt{dfs}(y_m)$ --- each $y_{i+1}$ is in $\texttt{graph}[y_i]$, and no intermediate $y_i$ equals $n-1$ (vertex $n-1$ occurs only at the end of a simple path). Different paths are different call chains, so each is recorded once.

\emph{Termination without marks.} The recursion depth is the length of the current path, and in a DAG every path has at most $n$ vertices. In a graph with a cycle the same code would recurse forever --- the marks of template 1 are what Problem~\ref{ssec:ch2-max-gold} needs. \hfill$\blacksquare$
\end{invariantbox}

\begin{alternativebox}[Approach 3: memoised path lists --- DP on the DAG]
\textbf{Idea.} The paths from $u$ to $n-1$ depend only on $u$, not on how $u$ was reached. So compute $P(u)$ = the list of paths $u \leadsto n-1$ once per vertex: $P(n-1) = [[n-1]]$, and $P(u)$ = ``$u$ in front of every path of $P(v)$'' for every successor $v$. This is DP on a DAG (\dpbook{12}{DP on DAGs \& Graphs}) with \emph{lists} as values.

\begin{lstlisting}[style=cpp]
#include <bits/stdc++.h>
using namespace std;

// Alternative: P(u) = all paths u -> n-1, computed once per vertex (memoised DFS = DP on the DAG).
class Solution {
    vector<vector<vector<int>>> memo;
    vector<char> done;
    const vector<vector<int>>& paths(const vector<vector<int>>& g, int u) {
        if (done[u]) return memo[u];
        vector<vector<int>> r;
        if (u == (int)g.size() - 1) r.push_back(vector<int>(1, u));
        else
            for (int v : g[u])
                for (auto& p : paths(g, v)) {          // every path v ... n-1, with u in front
                    vector<int> q(1, u);
                    q.insert(q.end(), p.begin(), p.end());
                    r.push_back(q);
                }
        done[u] = 1;
        memo[u] = r;
        return memo[u];
    }
public:
    vector<vector<int>> allPathsSourceTarget(vector<vector<int>>& graph) {
        memo.assign(graph.size(), {});
        done.assign(graph.size(), 0);
        return paths(graph, 0);
    }
};

int main() {
    vector<vector<int>> g = {{4, 3, 1}, {3, 2, 4}, {3}, {4}, {}};
    for (auto& p : Solution().allPathsSourceTarget(g)) {
        for (int x : p) cout << x << " ";
        cout << "/ ";
    }
    cout << "\n";                                      // 0 4 / 0 3 4 / 0 1 3 4 / 0 1 2 3 4 / 0 1 4 /
}
\end{lstlisting}

\begin{complexitybox}
\textbf{Time:} each vertex is expanded once (15 expansions on the complete DAG instead of $16\,384$ calls), but the lists are copied: building $P(u)$ costs the total length of $P(u)$, so the work is the sum of the sizes of \emph{all} lists, about twice the output on the complete DAG (measured: 9\,ms against 3\,ms for Approach~2, the copying dominates). \textbf{Space:} all lists $P(v)$ are kept at the same time. \textbf{Verdict:} \ac{}. Not faster here --- the output dominates --- but the same memo with \emph{numbers} instead of lists (count the paths, find the longest) is \VE{} where backtracking would be exponential.
\end{complexitybox}
\end{alternativebox}

\begin{tricksbox}[Approach 4: Implementation Refinements]
\begin{itemize}[leftmargin=*]
    \item \textbf{Prune dead ends.} Mark the vertices that can reach $n-1$ (a memoised DFS, or one search from $n-1$ on the reversed edges, Pattern~5) and never enter the others. On the DAG ``$0 \to 14$, plus $0 \to$ a complete DAG on $1..13$ that never reaches 14'' the plain backtracking makes $8193$ calls for one path; with the pruning it makes 2.
\begin{lstlisting}[style=cpp]
#include <bits/stdc++.h>
using namespace std;

// Backtracking that never enters a vertex from which n-1 cannot be reached.
class Solution {
    vector<vector<int>> res;
    vector<int> path;
    vector<int> good;                                  // -1 unknown, 0 no, 1 yes: v reaches n-1
    long long calls = 0;
    bool reaches(const vector<vector<int>>& g, int u) {
        if (good[u] != -1) return good[u];
        int ok = (u == (int)g.size() - 1);
        for (int v : g[u]) if (reaches(g, v)) ok = 1;  // memoised: every vertex is decided once
        good[u] = ok;
        return ok;
    }
    void dfs(const vector<vector<int>>& g, int u) {
        calls++;
        path.push_back(u);
        if (u == (int)g.size() - 1) res.push_back(path);
        else
            for (int v : g[u]) if (good[v] == 1) dfs(g, v);
        path.pop_back();
    }
public:
    vector<vector<int>> allPathsSourceTarget(vector<vector<int>>& graph) {
        good.assign(graph.size(), -1);
        res.clear(); calls = 0;
        if (reaches(graph, 0)) dfs(graph, 0);
        return res;
    }
    long long lastCalls() const { return calls; }
};

int main() {
    int n = 15;
    vector<vector<int>> g(n);
    for (int i = 1; i < n - 1; i++) {
        g[0].push_back(i);
        for (int j = i + 1; j < n - 1; j++) g[i].push_back(j);   // 1..13: complete, no way to 14
    }
    g[0].push_back(14);
    Solution s;
    size_t found = s.allPathsSourceTarget(g).size();
    cout << found << " " << s.lastCalls() << "\n";      // 1 2
}
\end{lstlisting}
    \item \textbf{One shared path vector.} Pushing and popping one \texttt{vector<int>} costs \bigO{1} per step; passing the path \emph{by value} into every call copies it at every node of the search tree.
    \item \textbf{Iterative form.} Not needed: the depth is at most $n = 15$. For the counting variant on big DAGs the memo is iterative over a topological order anyway.
    \item \textbf{Graphs with cycles.} Add \texttt{onPath} (template 1); without it the recursion never ends.
\end{itemize}
\end{tricksbox}

\subsubsection*{5. Trace}

\begin{tracebox}[{Example 2: graph = [[4,3,1],[3,2,4],[3],[4],[]]}]
\begin{center}
\begin{tikzpicture}[x=0.8cm, y=0.8cm]
    \tikzset{gv/.append style={minimum size=6.2mm, font=\scriptsize\bfseries}}
    \foreach \k/\t/\sa/\sb/\sc/\sd/\se in {
        0/{call 2: path [0,4] --- output}/gvseen/gv/gv/gv/gvcur,
        1/{call 4: path [0,3,4] --- output}/gvseen/gv/gv/gvseen/gvcur,
        2/{call 10: path [0,1,2,3,4] --- output}/gvseen/gvseen/gvseen/gvseen/gvcur}{
        \begin{scope}[xshift=\k*5.5cm]
            \node[\sa] (t0) at (0,0) {0}; \node[\sb] (t1) at (1.4,1.2) {1}; \node[\sc] (t2) at (2.9,1.2) {2};
            \node[\sd] (t3) at (2.9,-0.6) {3}; \node[\se] (t4) at (4.3,0) {4};
            \draw[gde] (t0) -- (t1); \draw[gde] (t0) -- (t3); \draw[gde] (t0) to[bend right=50] (t4);
            \draw[gde] (t1) -- (t2); \draw[gde] (t1) -- (t3); \draw[gde] (t1) -- (t4);
            \draw[gde] (t2) -- (t3); \draw[gde] (t3) -- (t4);
            \ifnum\k=0 \draw[gpath, -{Stealth[length=3mm]}] (t0) to[bend right=50] (t4); \fi
            \ifnum\k=1 \draw[gpath, -{Stealth[length=3mm]}] (t0) -- (t3); \draw[gpath, -{Stealth[length=3mm]}] (t3) -- (t4); \fi
            \ifnum\k=2 \foreach \a/\b in {0/1, 1/2, 2/3, 3/4} \draw[gpath, -{Stealth[length=3mm]}] (t\a) -- (t\b); \fi
            \node[font=\scriptsize\bfseries] at (2.15,2.0) {\t};
        \end{scope}
    }
\end{tikzpicture}
\end{center}
\begin{center}
\begin{tikzpicture}[tn/.style={rectangle, rounded corners=2pt, draw=darkgray!70, fill=white, font=\scriptsize\ttfamily, inner sep=2pt},
                    to/.style={tn, draw=accentorange, fill=answerfill, line width=0.9pt}]
    \node[tn] (c1) at (5,0) {1: 0};
    \node[to] (c2) at (0.6,-0.8) {2: 0 4};
    \node[tn] (c3) at (3.2,-0.8) {3: 0 3};
    \node[tn] (c5) at (8.4,-0.8) {5: 0 1};
    \node[to] (c4) at (3.2,-1.6) {4: 0 3 4};
    \node[tn] (c6) at (6.2,-1.6) {6: 0 1 3};
    \node[tn] (c8) at (8.4,-1.6) {8: 0 1 2};
    \node[to] (c11) at (10.8,-1.6) {11: 0 1 4};
    \node[to] (c7) at (6.2,-2.4) {7: 0 1 3 4};
    \node[tn] (c9) at (8.4,-2.4) {9: 0 1 2 3};
    \node[to] (c10) at (8.4,-3.2) {10: 0 1 2 3 4};
    \foreach \a/\b in {c1/c2, c1/c3, c1/c5, c3/c4, c5/c6, c5/c8, c5/c11, c6/c7, c8/c9, c9/c10}
        \draw[darkgray!80, -{Stealth[length=1.6mm]}] (\a) -- (\b);
    \node[note, text width=4.2cm, anchor=west] at (-0.4,-2.8) {the search tree: 11 calls, numbered in the order they happen; orange = an output};
\end{tikzpicture}
\end{center}
\begin{center}\small
\begin{tabular}{@{}rllll@{}}
\toprule
\textbf{call} & \textbf{enter} & \textbf{path after push} & \textbf{action} & \textbf{res} \\
\midrule
1 & 0 & \texttt{[0]} & try 4, 3, 1 & --- \\
2 & 4 & \texttt{[0,4]} & target: copy, pop & 1 path \\
3 & 3 & \texttt{[0,3]} & try 4 & \\
4 & 4 & \texttt{[0,3,4]} & target: copy, pop & 2 \\
5 & 1 & \texttt{[0,1]} & try 3, 2, 4 & \\
6 & 3 & \texttt{[0,1,3]} & try 4 & \\
7 & 4 & \texttt{[0,1,3,4]} & target & 3 \\
8 & 2 & \texttt{[0,1,2]} & try 3 & \\
9 & 3 & \texttt{[0,1,2,3]} & try 4 & \\
10 & 4 & \texttt{[0,1,2,3,4]} & target & 4 \\
11 & 4 & \texttt{[0,1,4]} & target & 5 \\
\bottomrule
\end{tabular}
\end{center}
Vertex 3 is entered three times (calls 3, 6, 9) and expands the same subtree $3 \to 4$ each time --- exactly what Approach~3 avoids: its memo computes $P(3) = [[3,4]]$ once and reuses it.
\end{tracebox}

\begin{diagrambox}
\begin{tikzpicture}[x=1.15cm, y=1cm]
    \node[gvsrc] (m0) at (0,0) {0}; \node[gv] (m1) at (1.4,1.2) {1}; \node[gv] (m2) at (2.9,1.2) {2};
    \node[gv] (m3) at (2.9,-0.6) {3}; \node[gvdst] (m4) at (4.3,0) {4};
    \draw[gde] (m0) -- (m1); \draw[gde] (m0) -- (m3); \draw[gde] (m0) to[bend right=50] (m4);
    \draw[gde] (m1) -- (m2); \draw[gde] (m1) -- (m3); \draw[gde] (m1) -- (m4);
    \draw[gde] (m2) -- (m3); \draw[gde] (m3) -- (m4);
    \node[font=\scriptsize\ttfamily, anchor=west, text=deepblue] at (4.7,0) {P(4) = [4]};
    \node[font=\scriptsize\ttfamily, anchor=north, text=deepblue] at (2.9,-0.95) {P(3) = [3 4]};
    \node[font=\scriptsize\ttfamily, anchor=south, text=deepblue] at (2.9,1.5) {P(2) = [2 3 4]};
    \node[font=\scriptsize\ttfamily, anchor=south east, text=deepblue, align=right] at (1.3,1.5) {P(1) = [1 3 4] [1 2 3 4] [1 4]};
    \node[font=\scriptsize\ttfamily, anchor=north east, text=deepblue, align=right] at (-0.2,-0.4) {P(0) = [0 4] [0 3 4]\\ \ \ \ \ [0 1 3 4] [0 1 2 3 4] [0 1 4]};
\end{tikzpicture}
\figcap{Approach 3 on Example 2: each vertex is expanded once (5 expansions instead of 11 calls); $P(0)$ comes out in the same order as the backtracking output.}
\end{diagrambox}

\subsubsection*{6. Test suite}

\begin{testbox}
\centering\small
\begin{tabular}{@{}p{5.6cm}p{3.6cm}p{5.2cm}@{}}
\toprule
\textbf{Input} & \textbf{Our output} & \textbf{Why} \\
\midrule
\texttt{[[1],[]]} & \texttt{[[0,1]]} & smallest graph, $n = 2$ \\
\texttt{[[],[]]} & \texttt{[]} & 0 has no edges: no path at all \\
\texttt{[[1],[],[1]]} & \texttt{[]} & $n-1 = 2$ is a source: unreachable \\
\texttt{[[2],[0],[]]} & \texttt{[[0,2]]} & vertex 0 is not the first vertex of the topological order \\
complete DAG, $n = 15$ & 8192 paths & the worst case: $2^{n-2}$ paths, 16\,384 calls \\
dead-end DAG above, $n = 15$ & \texttt{[[0,14]]} & 8193 calls plain, 2 calls with pruning \\
\bottomrule
\end{tabular}
\end{testbox}
All outputs come from running the programs above. Approaches 1--3 (and the pruned version) were also compared on $2000$ random DAGs with $n \le 10$ (random topological labels, random edge density, shuffled adjacency lists): the three sets of paths, sorted, were identical every time.

\subsubsection*{7. Common mistakes}

\begin{warningbox}[Mistake 1 --- a visited array copied from plain DFS]
\begin{lstlisting}[style=snippet]
if (!visited[v]) { visited[v] = true; dfs(g, v); }       // WRONG: each vertex used by one path only
\end{lstlisting}
On Example 2 this returns only \texttt{[[0,4]]}: vertex 4 is marked by the first path and can never end another one, and vertex 3 is never re-entered from 1. In a DAG drop the array; in a graph with cycles un-mark on exit.
\end{warningbox}

\begin{warningbox}[Mistake 2 --- recording the path without copying]
Storing a pointer or a reference to \texttt{path} instead of a copy gives a list of identical (finally empty) paths, because \texttt{path} keeps changing. \texttt{res.push\_back(path)} copies, which is what we want.
\end{warningbox}

\begin{warningbox}[Mistake 3 --- forgetting the pop on the target branch]
Returning early when $u = n-1$ (``\texttt{if (u == n-1) \{ res.push\_back(path); return; \}}'') skips the \texttt{pop\_back}; the vertex $n-1$ stays in the path and every later path contains it in the middle. Pop before returning, or keep a single exit as in Approach~2.
\end{warningbox}

\begin{warningbox}[Mistake 4 --- assuming 0 is a source]
Nothing says that 0 comes first in the topological order. \texttt{[[2],[0],[]]} has the edge $1 \to 0$; the backtracking does not care, but a brute force that orders vertices by number would.
\end{warningbox}

\subsubsection*{8. Takeaways}

\begin{itemize}[leftmargin=*]
    \item Backtracking is DFS plus a path stack; in a DAG the marks are unnecessary.
    \item When the output is exponential, the right goal is ``proportional to the output'': at most $2^{n-2}$ paths, and the search makes at most about twice as many calls.
    \item The future of a path in a DAG depends only on its last vertex --- memoise. With lists it gives the same output; with numbers it turns the exponential count into \VE.
    \item Prune what cannot reach the target before searching: a reverse search costs \VE{} and can save an exponential amount.
\end{itemize}

\begin{variantsbox}
\variant{Only the number of paths}{\varnone}
  {return how many paths lead from 0 to $n-1$, with $n$ up to $10^5$.}
  {$\text{cnt}(u)$ = number of paths $u \leadsto n-1$, depending only on $u$.}
  {$\text{cnt}(n-1) = 1$, $\text{cnt}(u) = \sum_{u \to v} \text{cnt}(v)$, memoised; \VE{}. On the complete DAG with $n = 15$ it returns 8192, the length of the list of Approach~2.}
\variant{The graph may have cycles}{\varin{https://atcoder.jp/contests/abc284/tasks/abc284_e}}
  {count the simple paths from vertex 1 in an undirected graph, capped at $10^6$.}
  {paths can now return to used vertices; the past matters.}
  {un-mark on exit (template 1) and stop at the cap: every call is one path (Exercise E4.4).}
\variant{Visit every vertex}{\varin{https://atcoder.jp/contests/abc054/tasks/abc054_c}}
  {count the paths from 1 that visit all $n \le 8$ vertices of an undirected graph.}
  {a path is a goal only when its length is $n$.}
  {backtracking with marks, answer counted at depth $n$ (Exercise E4.1).}
\variant{The longest path}{\varnone}
  {return one path from 0 to $n-1$ with the most vertices, $n$ up to $10^5$.}
  {$\text{best}(u)$ = the most vertices on a path $u \leadsto n-1$, plus the successor that achieves it.}
  {DP on the DAG, $\text{best}(u) = 1 + \max_{u \to v} \text{best}(v)$, then follow the stored successors; \VE. Both DP variants were checked against the lists of Approach~2 on 2000 random DAGs.}
\end{variantsbox}
